\scanpage{199}
par $\mathbf Q_\ell(\unsure{K|k})$ -- c'est particulièrement intéressant dans le cas $k$ un corps de nbs.

(iv) Si $x$ est un pt à valeurs \uline{dans $\mathbf C$}, il lui correspond
\[
\xi_{x,\mathbf Z}\in\mathbf T_\alpha(\mathbf Z)
\qquad \xi_{\unsure{\alpha|x}}(\mathbf Z)(F)=\text{réalisation de }F_x,
\]
\deleted{\unsure{Quand} -- De plus, le gr. $\mathbf Z/2\mathbf Z$ \unsure{opère sur} \inserted{$\xi_{x,\mathbf Z}$ et $\xi_{\unsure{\tau x},\mathbf Z}$ \unsure{sont can. isom.}} \illegible{3} ($\tau$ est la conj. complexe),}
et des isom. can.
\[
\left\{\begin{aligned}
\xi_{x,\mathbf Z}\otimes_{\mathbf Z}\mathbf Z_\ell&\simeq\xi_{x,\ell}\\
\xi_{x,\mathbf Z}\otimes_{\mathbf Z}\mathbf C&\simeq\xi_{x,\infty}
\end{aligned}\right.
\]

D'ailleurs $\xi_{x,\mathbf Z}$ est défini en termes d'un pt à val. dans un corps \inserted{top.} $K$ isom. à $\mathbf C$ (à isom. \unsure{can.} près) en d'autres termes on a un isom. canonique ``involutif''
\[
\xi_{x,\mathbf Z}\simeq\xi_{\tau x,\mathbf Z}
\]
où $\tau x$ est le pt complexe conjugué.

\deleted{\unsure{Donc}} On peut dire encore, si $x$ est localisé en $a\in X$, que \inserted{a) Si $k(a)$ est dense dans $\mathbf C$, alors} $\xi_{x,\mathbf Z}$ dépend de la donnée d'une topologie \inserted{archimédienne}, (ou valuation) sur $k(a)$, à complété $\mathbf C$.
\unsure{b) si $k(a)$ non dense dans}
\margin{198}
